\documentclass[12pt,english]{article}
\usepackage[T1]{fontenc}
\usepackage[latin9]{inputenc}
\usepackage[toc,page]{appendix}
\usepackage{geometry}
\geometry{verbose,tmargin=1.5in,bmargin=1.5in,lmargin=1in,rmargin=1in}
\usepackage{bm}
\usepackage{amstext}
\usepackage{comment}
\usepackage{amsmath}
\usepackage{hyperref}
\hypersetup{
     colorlinks = true,
     citecolor = navy,
     linkcolor = navy,
     urlcolor=navy
}
\usepackage[space]{grffile}
\usepackage{setspace}
\usepackage[normalem]{ulem}
\usepackage{xcolor}
\usepackage{soul}
\usepackage{dsfont}
\usepackage{refstyle}
\usepackage[capitalize,nameinlink]{cleveref}
\usepackage{setspace}
\usepackage{bbm}
\usepackage[tableposition=t]{caption} % ensures caption of longtables is normalsized.
\usepackage{mathpazo} % add Palatino font
\usepackage{longtable}
\usepackage{amssymb}
\usepackage{amsfonts}
\usepackage{amsmath}
%\usepackage[margin=1in]{geometry} % to set margins
\usepackage{scalefnt}
\usepackage{wrapfig}
\usepackage{graphicx}
%\usepackage{subfigure}
\usepackage{lscape}
\usepackage{pdflscape}
\usepackage{multirow}
\usepackage{array}
\usepackage{rotating}
\usepackage{bibunits}
\usepackage{fix-cm}
\usepackage{float}
\floatplacement{table}{H}
\usepackage{epsfig}
\usepackage{epstopdf}
\usepackage{refstyle}
\usepackage{booktabs}
\usepackage{tabularx}
\usepackage{array}

\newcolumntype{x}[1]{>{\raggedleft\arraybackslash}p{#1}}%
\newcolumntype{d}[1]{>{\centering\arraybackslash}p{#1}}%
\newcolumntype{y}[1]{>{\raggedright\arraybackslash}p{#1}}%
\usepackage{verbatim}
%\epstopdfsetup{suffix=}
\setcounter{MaxMatrixCols}{10}
\usepackage[longnamesfirst,authoryear]{natbib}
\newenvironment{proof}[1][Proof]{\noindent \textbf{#1.}}{\rule{0.5em}{0.5em}}
\setlength{\topmargin}{-0.80in} \setlength{\textheight}{9.50in}
\setlength{\evensidemargin}{.0in} \setlength{\oddsidemargin}{.0in}
\setlength{\textwidth}{6.85in}
\setlength{\voffset}{0.4cm}
\setlength{\hoffset}{-0.5cm}
\usepackage{setspace}
\usepackage{afterpage}
\usepackage[flushleft]{threeparttable}
\usepackage{tabulary}
%\usepackage[titletoc,title]{appendix}
\usepackage[bottom]{footmisc}
\usepackage{color}
	\definecolor{maroon}{RGB}{128,0,0}
	\definecolor{navy}{RGB}{0,0,165}
	\definecolor{lb}{RGB}{150,165,190}
	\definecolor{MyWhite}{RGB}{255, 255, 255}
\usepackage{colortbl}

\usepackage{caption}
\usepackage{subcaption}
\usepackage{placeins}
\usepackage{nicefrac}
%\usepackage{chngcntr}

\newcommand{\bit}{\begin{itemize}}
\newcommand{\eit}{\end{itemize}}
\usepackage{listings}
\def\cN{\mathcal{N}}
\newtheorem{theorem}{Theorem}
\newtheorem{remark}[theorem]{Remark}
\usepackage{mdframed}
\newmdtheoremenv{assumption}{Assumption}

\usepackage{enumitem}


\begin{document}


%\onehalfspacing


\title{Re-Re-thinking the ShapirIO approach\\ \emph{Using the Requirements Matrix to get what we want}}
%Beggar or Prosper-thy-Neighbour? \\
%Fiscal Spillovers Across Firms and Countries
\maketitle

\paragraph{Purpose:} Apologies all -- I was not sufficiently clear with my instructions about what to do with the Requirements Matrix that Emily found. This note is intended to correct that.

\section*{The theoretical framework}

We observe prices, consumption quantities and consumption expenditures for a variety of sectors $S$ and time periods $t$: $(p_{s,t},c_{s,t},FD_{s,t})_{s\in S,t\in \{1959Q1,...\}}$. We also observe a \emph{requirements} matrix $\Xi$ and know overall consumption expenditures by sector in 2017: $FD_{s,2017}$.

Before doing the original operations I specified, we need to compute two objects I use below: $\mathbold{\Omega}$ -- a matrix of bilateral cost shares where each element $\omega_{ij}$ denotes row industry $i$ sourcing a fraction $\omega_{ij}$ of its intermediates on goods from column industry $j$. That is rows will sum to 1. Note that this matrix is \emph{nominal} spending shares and not real in any way. So if the cellphone sector on average spent \$50 per cellphone on plastic and \$200 per cellphone on microchips (and there were no other intermediates), then the cellphone row will have an entry of 0.2 (50/(200+50)) under plastic and 0.8 (200/(200 + 50)) under microchips.

The other object is the share of \emph{sales} spent on intermediates. This is a row vector with elements $\gamma_s$ (let $\underline{\gamma}$ be the entire column vector). 

Note that the unit requirements matrix data we have (denoted $\Xi$) is given by:
$$
	\Xi \equiv (I - \mathbold{\Omega}'diag(\underline{\gamma}))^{-1}
$$

Now this means we can construct our $\mathbold{\Omega}$ and $\underline{\gamma}$ as:

\begin{align*}
	\underline{\gamma} &= (I-\Xi^{-1})1_S &\text{i.e. rowsum of }(I-\Xi^{-1})\\
	\Omega &= [(I-\Xi^{-1})diag(1/\underline{\gamma})]'& \substack{\text{i.e. divide elements of the requirements}\\\text{matrix by their rowsum (and transpose)}}
\end{align*}

Now we proceed with the regular procedure using these calculations of $\Omega$ and $\underline{\gamma}$. Alternatively, wherever you see $(I-\Omega'\underline{\gamma})^{-1}$ you can replace it with $\Xi$. \emph{But don't forget for calculating intermediate prices to use the $\Omega$ defined in this note!}

Next follow the old procedure below:

\section*{Old Procedure}


We can easily calculate sales in 2017 as:
\[
	Sales_s = \sum_{j}IO_{j,s}+FD_s
\]

But we would like a method to calculate sales in all time periods.

To do this, let's assume that we have the same I-O matrix \emph{in shares} as we do in 2017. With that assumption we can now use our time series of consumption expenditures $FD_{s,t}$ to get:

\[
	Sales_{s,t} = (I-\mathbold{\Omega}'diag(\underline{\gamma}))^{-1}FD_{s,t} 
\]
where FD can be the expenditure series (in USD) or it can be the product of the price and quantity indexes $\widetilde{FD}_{s,t}=p_{s,t}\times c_{s,t}$.

While we have \emph{output} prices for all goods it would be nice to  have a sense of the price of value added in each sector (i.e. removing the price of intermediates). So let's see if we can calculate this. Under Cobb-Douglas production we have:

\[
	P_s^{int} = \exp^{\sum_{j}\Omega_{i,j}\log(P_{j,t})}
\]

More generally, we have:

\[
	P_{s,t}^{int} = \left[\sum_{j}\Omega_{i,j}P_{j,t}^{1-\theta}\right]^{\frac{1}{1-\theta}}
\]
where $\theta$ is the elasticity of substitution between intermediate varieties. Ideally this would be set to a low value.

Then if we assume a similar structure between value added and intermediates (and assume a gross markup of 1) then we get one of:

\begin{align*}
	mc_{s,t}&=\exp^{(1-\gamma_s)\log(P_{s,t}^{VA})+\frac{\gamma_s}{1-\mu_{s,t}}\log(P_{s,t}^{int})} & \text{Cobb-Douglas Production}\\
mc_{s,t}&=\left[(1-\gamma_s)(P_{s,t}^{VA})^{1-\epsilon}+\gamma_s(P_{s,t}^{int})^{1-\epsilon}\right]^{\frac{1}{1-\epsilon}} & \text{Constant Elasticity of Substitution Production}	
\end{align*}
$P_{s,t}^{VA}$ is the price of value added in sector $s$ (this is the cost of labor plus land and capital required to produce a unit of output). 

Again assuming that markups are 1, we can calculate a time-series for $P_{s,t}^{VA}$ that is correct up to a (sector-specific) scale factor:

\begin{align*}
	mc_{s,t}&=\frac{P_{s,t}}{\mu_s}\\
	\Rightarrow P_{s,t}^{VA} &= \exp^{\frac{1}{1-\gamma_s}\left[\log(P_{s,t})-\gamma_s\log(P_{s,t}^{int})\right]}&\text{Cobb-Douglas}\\
	\Rightarrow P_{s,t}^{VA} &= \left(\frac{1}{1-\gamma_s}\left[\left(P_{s,t}\right)^{1-\epsilon}-\gamma_s(P_{s,t}^{int})^{1-\epsilon}\right]\right)^{\frac{1}{1-\epsilon}}&\text{CES}
\end{align*}

\section{A new procedure to follow:}

Instead of running the VAR first and then doing the I-O adjustments, let's do the I-O adjustment first and then do the VAR. So this would entail:

\begin{enumerate}
	\item Find $sales_{s,t}$ using the formulation above.
	\item Find $P_{s,t}^{VA}$ as well.
	\item Compare $P_{s,t}^{VA}$ to $P_{s,t}$ using the bar graph approach we already used (pick 1 quarter every decade to do this).
	\item If 3. checks out, then run two sets of VARs: one matching Shapiro with $(p_{s,t},q_{s,t})$ and one with $(p_{s,t}^{VA},sales_{s,t})$. If you prefer, feel free to run a VAR with some notion of real production: $y_{s,t}=\frac{sales_{s,t}}{p_{s,t}}$ noting that we use the output price data $p_{s,t}$ here.
	\item Classify shocks (sector-by-sector and time period by time period) as mostly demand or mostly supply using the Shapiro classification.
	\item Also classify the mostly demand/mostly supply as positive or negative
	\item Compare the IO vs. Shapiro classifications.
\end{enumerate}


\end{document}